\section{Produzenten}

\begin{frame}{Produzenten: Thermische Kraftwerke}
    \mycomment{
    Bei thermischen Kraftwerke treiben Brennstoffe Turbinen und diese wiederum Generatoren an. Sie brauchen Kühlung, und sie erzeugen immer auch Wärme, die entweder als Abwärme verloren geht oder (zum Teil) genutzt wird, beispielsweise als Fernwärme (Kraft-Wärme-Kopplung).

    Diese Turbinen erzeugen direkt die benötigte Frequenz 50 Hertz.

    Es ist möglich, mit der entstehenden Wärme sofort einen weiteren Generator zu betreiben, der ebenfalls Strom erzeugt (GuD-Kraftwerk, Gas- und Dampfturbinenprozess).

    Der produzierte Strom kann natürlich auch andernorts zur Heizung verwendet werden, aber auch zur Kühlung (Klimaanlagen). 
    }
    \showmytikz[1]{../shared/tikz/tikz-energie-kws}
\end{frame}

\begin{frame}{Kinetische Energie}
    \mycomment{
    Fossile (und organische) Brennstoffe können direkt genutzt werden, um Verbrennermotoren zu betreiben. Strom kann genutzt werden, um Elektromotoren zu betreiben. (Eine weitere Nutzung ist natürlich Licht und Elektronik, die allerdings kaum \enquote{zur Last} fällt.)
    }
    \showmytikz[1,2]{../shared/tikz/tikz-energie-kws}
\end{frame}

\begin{frame}{Produzenten: Erneuerbare Energiequellen}
    \mycomment{
    Erneuerbare erzeugen Strom durch Direktumwandlung. Wind und Wasser erzeugen mit Turbinen Wechselstrom; bei Wind wird dann ein Gleichrichter benötigt. Wasserkraftwerke können mitunter direkt 50 Hertz liefern.     
    Solarzellen erzeugen Gleichstrom, der umgewandelt werden muss. 
    
    }
    \showmytikz[1,2,3]{../shared/tikz/tikz-energie-kws}
\end{frame}

\begin{frame}{Neue Entwicklungen: Sektorkopplung}
    \mycomment{
    Strom kann auch genutzt werden, um Wärme zu erzeugen. Auch kann Sonnenergie in Form von Wärme gespeichert werden. Jede solche Umwandlung ist allerdings mit Verlusten verbunden.
    
    }
    \showmytikz[1,2,3,4]{../shared/tikz/tikz-energie-kws}
\end{frame}

\begin{frame}{Neue Entwicklungen: Prosumer}
    \mycomment{
    Die Energiemärkte sind im Wandel, und das nicht nur wegen steigender Anteile der Erneuerbaren. Die Lasten verändern sich und werden zunehmend \enquote{intelligent}. Die Netze müssen an das veränderte Nutzerverhalten angepasst werden - ebenso wie Verträge.
    }
    \mytext{
    \begin{itemize}
        \item die klassische Unterscheidung zwischen Leistung und Last (Produzenten und Verbrauchern) verschwimmt
        \item Einspeisung durch z.B. Balkonkraftwerke, Rückeinspeisung von Electric Vehicles
        \item Energiegemeinschaften, z.B. ganze Dörfer, mit Lasten und Erzeugung
        \item Smart Meter: ggf. Möglichkeiten zur Abregelung
        \item Variable vs. fixeStromtarife
    \end{itemize}
    }
\end{frame}